\chapter{采样离散控制与滤波的实际影响}
\label{ch:digital}

连续控制方程描述每一瞬间，程序只在一串离散时刻执行。每个周期依次获得数据、计算、发送命令，驱动器再在两个更新之间保持上一条命令。这些步骤改变了闭环，不能把连续公式逐字翻成代码后就认为性能不变。本章以 $200\,\mathrm{Hz}$、即 $T_s=5\,\mathrm{ms}$ 的外部控制周期为主要教学例子；它只代表一种采样设置，不代表任何机械臂的固有带宽。

\section{先分清采样量化保持与延迟}

采样是只在 $t_k=kT_s$ 读信号；量化是把连续幅值舍入到有限分辨率；保持是两次更新之间维持某个输出；延迟是数据所代表的时刻与输入真正生效的时刻不同。这四件事可以同时存在，修正其中一项不会自动消除其他三项。

例如一条关节状态在 $t=1.000\,\mathrm{s}$ 被测量，主机在 $1.003\,\mathrm{s}$ 才收到，即使计算只用了 $0.2\,\mathrm{ms}$，控制器使用的状态也已经旧了 $3.2\,\mathrm{ms}$。仅打印计算函数耗时不足以描述闭环延迟，必须记录测量时间、接收时间、命令发送时间和实际生效时序。

\subsection{采样后高频会伪装成低频}

对于连续正弦 $x(t)=\sin(2\pi ft)$，采样后是 $x[k]=\sin(2\pi f k/f_s)$。把 $f$ 换成 $f+mf_s$，$m$ 为整数，采样结果不变，因为多出来的是 $2\pi mk$。这就是混叠的来源，不是软件计算误差。

\begin{examplebox}[算例 13A 170 Hz 为什么会看起来像 30 Hz]
取 $f_s=200\,\mathrm{Hz}$：
\[
\sin(2\pi170k/200)
=\sin(2\pi k-2\pi30k/200)
=-\sin(2\pi30k/200).
\]
同一串采样点既可由 170 Hz 信号产生，也可由负相位的 30 Hz 信号产生。软件拿到采样点后已经无法区分两者。应在采样前做合适的抗混叠处理，或先以更高频率采样并滤波后再降采样；不是在 200 Hz 数据上随意加一个低通就能恢复丢失的信息。
\end{examplebox}

\begin{figure}[htbp]
\centering\includegraphics[width=.95\linewidth]{figures/cf_sampling.pdf}
\caption{左图的高频信号与低频信号通过相同采样点；右图展示零阶保持的阶梯输出。二者分别是信息混叠和输出保持，不能混为一个问题。}
\end{figure}

采样定理要求信号有效带限且 $f_s>2f_{\max}$。实际还要给抗混叠滤波器的过渡带留空间。控制设计则比信号重建更严格，因为采样虽能记录某个频率，闭环仍可能没有足够相位余量去控制该频率。

\section{从连续模型精确得到一拍更新}

对 $\dot x=Ax+Bu$，在区间 $[kT_s,(k+1)T_s)$ 内保持 $u=u_k$，解微分方程得
\begin{equation}
x_{k+1}=A_dx_k+B_du_k,
\qquad A_d=e^{AT_s},\qquad
B_d=\int_0^{T_s}e^{A\xi}B\,d\xi.
\label{eq:cf_zoh_exact}
\end{equation}
这叫零阶保持等效离散化。它对区间内恒定输入在采样点的状态是精确的；不意味着区间内轨迹变成直线，也不包含额外通信延迟。

\subsection{手算教学关节的精确离散模型}

速度方程为 $\dot v=-(b/J_r)v+u/J_r$。一拍后
\[
v_{k+1}=a v_k+\frac{1-a}{b}u_k,
\qquad a=e^{-bT_s/J_r}.
\]
位置由区间内速度积分：
\begin{align}
q_{k+1}
&=q_k+\int_0^{T_s}
\left[e^{-b\xi/J_r}v_k+\frac{1-e^{-b\xi/J_r}}b u_k\right]d\xi\\
&=q_k+\frac{J_r}{b}(1-a)v_k
+\left[\frac{T_s}{b}-\frac{J_r}{b^2}(1-a)\right]u_k.
\end{align}
对 $J_r=0.1$、$b=0.2$、$T_s=0.005$，有 $a=e^{-0.01}=0.9900498$，因此
\begin{equation}
A_d=\begin{bmatrix}1&0.00497508\\0&0.99004983\end{bmatrix},
\qquad B_d=\begin{bmatrix}0.000124584\\0.04975083\end{bmatrix}.
\end{equation}

\begin{examplebox}[算例 13B 相同一拍用三种积分方法会怎样]
从 $q_0=v_0=0$ 开始，保持 $u=1\,\mathrm{N\,m}$ 一拍。精确结果是
\[
v_1=0.0497508\,\mathrm{rad/s},\qquad
q_1=0.000124584\,\mathrm{rad}.
\]
前向欧拉先用旧状态算导数，给 $v_1=0.05$、$q_1=0$；半隐式欧拉先更新速度再更新位置，给 $v_1=0.05$、$q_1=0.00025$。两种数值法都在近似同一个连续对象，但一拍的位置分别偏小和偏大。控制周期变小能减小误差，积分方法本身仍需选清楚。
\end{examplebox}

\section{连续稳定不等于离散实现稳定}

离散线性系统 $x_{k+1}=A_dx_k$ 的渐近稳定条件是全部特征值严格位于单位圆内，即 $|z_i|<1$。若用精确映射 $z=e^{sT_s}$，连续极点 $\Re(s)<0$ 对应 $|z|<1$。但前向欧拉把 $\dot x=-ax$ 变成 $x_{k+1}=(1-aT_s)x_k$，于是要求
\[
|1-aT_s|<1\quad\Longleftrightarrow\quad0<aT_s<2.
\]

\begin{examplebox}[算例 13C 一个本来稳定的衰减被欧拉算成发散]
连续系统 $\dot x=-500x$ 显然指数衰减。取 $T_s=0.005\,\mathrm{s}$，精确离散极点为 $e^{-2.5}=0.082085$，前向欧拉极点却是 $1-2.5=-1.5$。从 $x_0=1$ 开始，欧拉序列是 $1,-1.5,2.25,-3.375,\ldots$，符号交替且幅值越来越大。物理模型并未变得不稳定，失稳来自数值实现。
\end{examplebox}

\begin{figure}[htbp]
\centering\includegraphics[width=.95\linewidth]{figures/cf_discrete_stability.pdf}
\caption{连续稳定系统的精确更新与不合适的前向欧拉更新。单位圆判断针对离散极点，不能继续只看连续极点在不在左半平面。}
\end{figure}

后向欧拉的替换为 $s\leftarrow(z-1)/(T_sz)$；Tustin 双线性变换为
\[
s\leftarrow\frac{2}{T_s}\frac{z-1}{z+1}.
\]
它们对连续左半平面的映射具有更好的稳定性性质，但仍会改变频率和瞬态。后向欧拉可能把部分连续不稳定动态数值上压住，因而“看起来不发散”不能用来确认真实系统稳定。Tustin 对频率有扭曲，设计窄带陷波时通常需要预畸变。

\section{给相位建立一张预算表}

纯延迟 $T_d$ 的传递函数是 $e^{-sT_d}$。令 $s=j\omega$，幅值为 1，相位为 $-\omega T_d$。零阶保持在相对较低频率处可近似贡献半拍延迟；如果还要等下一拍输出，则常见的一种时序近似为总延迟 $1.5T_s$。这不是所有控制器的固定常数，实际值要根据时间戳和执行结构确定。

\begin{examplebox}[算例 13D 200 Hz 更新不代表能闭环控制到 200 Hz]
取 $T_s=5\,\mathrm{ms}$，假设一拍计算等待加半拍保持，总等效延迟 $7.5\,\mathrm{ms}$。相位损失为
\[
\Delta\phi=360^\circ fT_d.
\]
在 5、10、20 Hz 处分别为 $13.5^\circ$、$27^\circ$、$54^\circ$。若设计前已有 $65^\circ$ 相位裕度，想把穿越频率放在 20 Hz，单是这一项就只剩约 $11^\circ$。

再考虑 20 Hz 截止的一阶低通，在 10 Hz 处损失 $\arctan(10/20)=26.565^\circ$。若穿越仍大约在 10 Hz，延迟与滤波合计约 $53.6^\circ$。真实穿越点也可能因滤波而改变，因此最终必须用完整开环重新求穿越，而不是只做度数相减。
\end{examplebox}

“每秒算 200 次”只说明更新率。可用带宽还由力矩内环、机械柔性、测量滤波、通信、计算和保持共同决定。阻尼增益增大后看见抖动，不应立刻认定“阻尼不够”；速度噪声和相位滞后可能已经让名义阻尼变成不合适的能量输入。

\section{低通滤波的公式怎样变成代码}

一阶连续低通 $H(s)=a/(s+a)$ 对应 $\dot y=a(x-y)$，其中 $a=2\pi f_c$。若输入在一拍内保持为 $x_k$，精确递推为
\[
y_{k+1}=\alpha y_k+(1-\alpha)x_k,
\qquad\alpha=e^{-aT_s}.
\]
这里的下标明确表示：当前输入 $x_k$ 作用到下一拍状态 $y_{k+1}$。某些实现收到当前采样后立即更新并把新值称为 $y_k$，只是索引约定不同，但与其他信号同步时必须知道它代表哪个时刻。

后向欧拉常见的“收到当前输入立即更新”形式为
\[
y_k=\frac{1}{1+aT_s}y_{k-1}
+\frac{aT_s}{1+aT_s}x_k.
\]
Tustin 则给出
\[
y_k=\frac{2-aT_s}{2+aT_s}y_{k-1}
+\frac{aT_s}{2+aT_s}(x_k+x_{k-1}).
\]
三者输入时序和频率响应并不完全相同，不能只把不同的系数都塞进同一条递推式然后称作完全等价。

\begin{examplebox}[算例 13E 计算 20 Hz 一阶低通系数]
取 $f_c=20\,\mathrm{Hz}$、$T_s=0.001\,\mathrm{s}$，得到 $aT_s=0.125664$。精确更新中新输入权重为 $1-e^{-0.125664}=0.118089$；后向欧拉的新输入权重为 $0.125664/1.125664=0.111635$；Tustin 的两个输入系数均为 $0.0591174$，旧输出系数为 $0.881765$。

若使用后向欧拉，旧输出为 1，当前输入为 2，则新输出为 $1+0.111635(2-1)=1.111635$。参数越靠近采样能力的上限，不同离散方法的差异越不能忽略。
\end{examplebox}

\subsection{更高阶不自动代表更适合闭环}

Butterworth 低通追求平坦通带；Chebyshev 允许纹波换取陡峭过渡；Bessel 侧重通带内较平坦的群延迟；椭圆滤波器同时允许通带和阻带纹波以获得更陡过渡。这些是设计目标，不是优劣排名。闭环关心的是噪声衰减与相位一起变化后的结果。

高阶 IIR 常拆成二阶节，即 SOS。对一个归一化二阶节
\[
H(z)=\frac{b_0+b_1z^{-1}+b_2z^{-2}}{1+a_1z^{-1}+a_2z^{-2}},
\]
对应差分式为
\[
y_k=b_0x_k+b_1x_{k-1}+b_2x_{k-2}-a_1y_{k-1}-a_2y_{k-2}.
\]
公式右边减号来自把分母乘到左边后移项。把高阶多项式拆成二阶节通常能改善浮点数值条件，但稳定性仍需检查每节极点、系数精度和状态初始化，不能只因为用了 SOS 就保证不会出问题。

\begin{figure}[htbp]
\centering\includegraphics[width=.98\linewidth]{figures/cf_filters.pdf}
\caption{低通同时改变幅值和相位；双向滤波把幅值平方；陷波只针对一个频带。中图采用连续 Butterworth 幅频公式展示原理，实际数字滤波的精确频响应以所用系数计算。}
\end{figure}

\subsection{离线零相位不能搬进实时回路}

正向滤波再反向滤波需要未来数据，是非因果处理。对实系数滤波器，理想稳态频率响应变成
\[
H_{\rm fb}(e^{j\Omega})=H(e^{j\Omega})H(e^{-j\Omega})
=|H(e^{j\Omega})|^2.
\]
相位被抵消，幅值却不保持原样。单向截止处幅值为 $1/\sqrt2$，双向后变成 $1/2$，对应约 $-6.02\,\mathrm{dB}$，不是 $-3.01\,\mathrm{dB}$。边界填充和初值处理还会产生端点效应。

对四阶模拟 Butterworth 原型，单向幅值为 $[1+(f/f_c)^8]^{-1/2}$，双向后为 $[1+(f/f_c)^8]^{-1}$。若名义 $f_c=8\,\mathrm{Hz}$，双向结果在 8 Hz 为 0.5。实际离散滤波要考虑频率映射。第9章的辨识数据处理中还必须注意：同一个线性滤波算子应施加到回归等式两边；仅仅分别把 $q$ 和 $\tau$ 平滑一下，并不能保证非线性构造的回归矩阵仍与力矩完全同步。

\subsection{移动平均和局部多项式}

$N$ 点滑动平均是有限脉冲响应 FIR 滤波器。对称线性相位 FIR 的群延迟为 $(N-1)T_s/2$。例如 31 点、200 Hz 采样会带来 $75\,\mathrm{ms}$ 延迟；在 5 Hz 处对应 $135^\circ$ 相位。这种处理可能适合离线看曲线，却很难直接放进要求较快响应的反馈路径。

Savitzky--Golay 方法在小窗口内拟合多项式，再由拟合多项式取值或求导。五点二次拟合的一阶和二阶导数为
\begin{align}
\dot y_k&\approx\frac{-2y_{k-2}-y_{k-1}+y_{k+1}+2y_{k+2}}{10T_s},\\
\ddot y_k&\approx\frac{2y_{k-2}-y_{k-1}-2y_k-y_{k+1}+2y_{k+2}}{7T_s^2}.
\end{align}
它不是不知道未来也能算出的神奇微分器；上述居中窗口需要 $k+1,k+2$ 两个未来样本。若实时实现为延迟输出，须计入两拍以上的等待。

\begin{examplebox}[算例 13F 用五个位置样本计算速度和加速度]
取 $q(t)=t^2$，$T_s=0.1\,\mathrm{s}$，在中心时刻 $t=0.2$ 周围的五个样本为 $[0,0.01,0.04,0.09,0.16]$。代入一阶公式：
\[
\dot q=\frac{0-0.01+0.09+0.32}{1}=0.4.
\]
代入二阶公式：
\[
\ddot q=\frac{0-0.01-0.08-0.09+0.32}{0.07}=2.
\]
与解析导数 $2t=0.4$、$2$ 一致，因为窗口中真实曲线恰好是二次多项式。遇到更复杂的高频运动或噪声时，它不再精确，窗口长度仍存在分辨率与平滑程度的权衡。
\end{examplebox}

\section{量化噪声为什么会在求导后放大}

若一圈有 $N_c$ 个有效计数，位置分辨率为 $\Delta q=2\pi/N_c$。这里应使用解码后的实际计数，不要把编码器“线数”和四倍频后的“计数数”混为一谈。直接差分速度为
\[
\hat v_k=\frac{q_k-q_{k-1}}{T_s}.
\]
一个计数跳变对应 $\Delta q/T_s$ 的速度台阶。若每次量化误差独立、均匀分布在 $[-\Delta q/2,\Delta q/2]$，单个位置误差方差为 $\Delta q^2/12$，两次差值的方差相加，因此
\[
\sigma_v=\frac{\Delta q}{\sqrt6\,T_s}.
\]
这个独立均匀假设只用于估算；真实低速编码器误差常呈周期性和相关性。

\begin{examplebox}[算例 13G 17 位编码器在 200 Hz 下的速度台阶]
$N_c=2^{17}=131072$，所以 $\Delta q=4.79369\times10^{-5}\,\mathrm{rad}$。在 $T_s=0.005\,\mathrm{s}$ 下，一个计数形成 $0.0095874\,\mathrm{rad/s}$ 的速度台阶。独立量化误差模型给出 $\sigma_v\approx0.003914\,\mathrm{rad/s}$。

若把采样周期减到 1 ms 而不改变差分方法，一个计数的速度台阶扩大到 $0.0479369\,\mathrm{rad/s}$。因此更高采样率不自动让简单差分更平滑；它需要与估计窗口、模型和滤波一起设计。
\end{examplebox}

\section{一拍控制更新必须明确状态顺序}

下列是带反算抗饱和和测量微分滤波的结构性伪代码。它用本书的 $e=q-q_d$，积分变量直接保存力矩贡献。这里没有机械臂接口，目的是看清更新顺序。
\begin{lstlisting}[language=C++,basicstyle=\ttfamily\small]
// q_prev must be initialized from the first measurement.
error = q - q_des;
velocity_raw = (q - q_prev) / dt;
velocity_f += alpha * (velocity_raw - velocity_f);
u_raw = -kp * error - kd * velocity_f + i_term + u_ff;
u_sent = clamp(u_raw, u_min, u_max);
i_term += dt * (-ki * error + kaw * (u_sent - u_raw));
q_prev = q;
\end{lstlisting}

积分更新在输出之后或之前均可设计，但推导、代码和日志必须一致。发生丢帧时，不应悄悄用固定 $T_s$ 对实际间隔不同的两个样本求导。使用实测间隔可以修正差分比例，却不能消除不规则采样本身造成的闭环变化。

增量式 PID 把 $u_k-u_{k-1}$ 作为计算结果，但积分信息仍藏在历史输出中。它不会天然免疫积分饱和，也不会自动实现无扰切换。切换模式时应初始化内部状态，使新控制器当前输出与切换前一致；对式\eqref{eq:cf_aw}，可令
\[
i_0=u_{\rm old}+K_pe+K_d\dot q-u_{\rm ff}.
\]
若新的输出限制更窄，则还要考虑限幅导致的实际变化。

\section{计算练习与逐步解答}

\textbf{练习 1}\quad 控制周期为 2 ms，一拍计算延迟加零阶保持近似为 $1.5T_s$。在 20 Hz 处损失多少相位？

\textbf{解}\quad 总延迟为 $0.003\,\mathrm{s}$，相位损失为 $360^\circ\times20\times0.003=21.6^\circ$。如果额外通信还有 1 ms，新增 $7.2^\circ$，合计 $28.8^\circ$。

\textbf{练习 2}\quad $\dot x=-80x$ 用前向欧拉，求允许的采样周期范围。取 10 ms 和 30 ms 时分别会怎样？

\textbf{解}\quad 稳定要求 $0<80T_s<2$，即 $0<T_s<0.025\,\mathrm{s}$。10 ms 时离散极点为 0.2，单调衰减；30 ms 时极点为 $-1.4$，交替发散。精确离散在 30 ms 时的极点为 $e^{-2.4}=0.090718$，仍然稳定。

\textbf{练习 3}\quad 采样率 500 Hz，使用 21 点对称 FIR。求群延迟，在 10 Hz 处相当于多少相位？

\textbf{解}\quad 延迟 $(21-1)/(2\times500)=0.02\,\mathrm{s}$，相位为 $-360^\circ\times10\times0.02=-72^\circ$。这是固定群延迟，不是 CPU 算得慢产生的延迟；换更快 CPU 无法消除它。

\textbf{练习 4}\quad 单向滤波器在某频率幅值为 0.8，相位为 $-40^\circ$。理想双向滤波后的幅值与相位是什么？

\textbf{解}\quad 幅值为 $0.8^2=0.64$，相位为 $-40^\circ+40^\circ=0$。幅值对应 $20\log_{10}(0.64)=-3.876\,\mathrm{dB}$。零相位并不表示没有改变幅值，也不意味着可实时执行。

\section*{本章与原稿的对应}
本章整合 [A，第18--22页] 与 [B，第67--70、78--83页]。补足精确 ZOH 数值、索引时序、混叠推导、量化噪声方差及五点求导计算，修正了“增量式自然无扰”“双向滤波仍是原来的截止定义”等容易误用的表述。
